\scanpage{80}
% archive foliation: 79
\margin{$\overline{\mathrm{I}}$ et $\overline{\mathrm{II}}$}

Rappels sur \deleted{sur}la théorie de Dieudonné (sur corps parfait k).

\margin{\uuline{Serre}, \uline{Corps Locaux}\\
\uuline{Demazure--Gabriel} Schémas en Groupes}

1. Rappels sur vecteurs de Witt.

\margin{$p$ est premier fixé dans tout le cours}

$\underline{W}_n\inserted{(n\geq1)}$, $\underline{W}$ schémas en anneaux sur $\underline{Z}$, $\underline{W}=\varprojlim\underline{W}_n$, avec
$\underline{W}_n=\underline{E}^{n}$ (absolu) comme schéma, $\underline{W}=\underline{E}^{\underline{N}}$ comme schéma, morphismes de transition
\[
R_n:\underline{W}_{n+1}\longrightarrow\underline{W}_n,\quad
R_n(x_1,\ldots,x_{n+1})=(x_1,\ldots,x_n)
\]
Structures additive et multiplicative caractérisées\deleted{p} par la condition que les \deleted{homo}morphismes de schémas
\[
\deleted{\unsure{W}}\qquad
\inserted{\Phi=(\Phi_1,\ldots,\Phi_n):}\quad
\underline{W}_n\longrightarrow\underline{O}^{\inserted{n}}\quad
(=\underline{E}^{\inserted{n}}\text{ comme schéma})
\]
donnés par'
\[
\phi_i(x_1,\ldots,x_n)=\phi_i(x_1,\ldots,x_i)
=x_1^{p^{i-1}}+px_2^{p^{i-2}}+\cdots+p^{i-1}x_i
\]
soi\deleted{en}t \deleted{des}\inserted{\unsure{un}} homomorphisme\deleted{s} d'\deleted{a}\inserted{A}nneaux.

\deleted{Homo}M\uline{orphismes remarquables:} \inserted{Morphisme \uline{additif}}
\[
V_n=V:\underline{W}_n
\xrightarrow[\deleted{\unsure{T}}]{\deleted{\unsure{T}}}\underline{W}_n,\qquad
V(x_1,\ldots,x_n)=(0,x_1,\ldots,x_{n-1})
\]
\deleted{donnés par} d'où par passage à la limite
\[
\deleted{\unsure{R}(x_1,\unsure{\ldots})}\qquad
V:\underline{W}\longrightarrow\underline{W},
\]
puis \deleted{\unsure{morph.}}
\[
\inserted{\text{morphisme d'Anneaux}}\quad
R_n\text{ ou }R:\underline{W}_{n+1}\longrightarrow\underline{W}_n,
\quad R(x_1,\ldots,x_{n+1})=(x_1,\ldots,x_n)
\]
\[
\inserted{\text{morphisme add.}}\quad
\unsure{\Pi|T}_n\text{ ou }\unsure{\Pi|T}:\underline{W}_n\longrightarrow\underline{W}_{n+1},
\quad\unsure{\Pi|T}(x_1,\ldots,x_n)=(0,x_1,,x,x_n)
\]
avec
\[
R_nT_n=V_n,\quad T_nR_n=V_{n+1}.
\]
Enf\unsure{in}\deleted{\unsure{car. p}}, on a
\[
F\text{ ou }F_n:\underline{W}_n\longrightarrow\underline{W}_n,\quad
F(x_1,\ldots,x_n)=(x_1^p,\ldots,x_n^p),
\]
d'où par passage à la limite
\[
F:\underline{W}\longrightarrow\underline{W}
\]
\uline{En caractéristique} $p$
\[
F_{\underline{F}_p}:\underline{W}_{\underline{F}_p}\longrightarrow\underline{W}_{\underline{F}_p}
\qquad(\text{et variante avec }\underline{W}_{n\underline{F}_p})
\]
